
\documentclass{article}
\usepackage{colortbl}
\usepackage{makecell}
\usepackage{multirow}
\usepackage{supertabular}

\begin{document}

\newcounter{utterance}

\centering \large Interaction Transcript for game `clean\_up', experiment `1\_medium\_5obj\_es', episode 1 with gpt{-}5.4{-}azure.
\vspace{24pt}

{ \footnotesize  \setcounter{utterance}{1}
\setlength{\tabcolsep}{0pt}
\begin{supertabular}{c@{$\;$}|p{.15\linewidth}@{}p{.15\linewidth}p{.15\linewidth}p{.15\linewidth}p{.15\linewidth}p{.15\linewidth}}
    \# & $\;$A & \multicolumn{4}{c}{Game Master} & $\;\:$B\\
    \hline

    \theutterance \stepcounter{utterance}  
    & & \multicolumn{4}{p{0.6\linewidth}}{
        \cellcolor[rgb]{0.9,0.9,0.9}{
            \makecell[{{p{\linewidth}}}]{
                \texttt{\tiny{[P1$\langle$GM]}}
                \texttt{Soy tu maestro de juego, y estás jugando un juego colaborativo con la siguiente cuadrícula como tablero de juego:} \\
\\ 
\texttt{```} \\
\texttt{ 1234567} \\
\texttt{╔══╤══╤═╗} \\
\texttt{║◌◌│T◌│◌║ 1} \\
\texttt{╟──┼──┴─╢ 2} \\
\texttt{║WI│◌◌◌◌║ 3} \\
\texttt{╟──┘◌◌◌◌║ 4} \\
\texttt{║◌◌◌H◌◌C║ 5} \\
\texttt{║◌◌┌─┬──╢ 6} \\
\texttt{║◌◌│◌│◌◌║ 7} \\
\texttt{╚══╧═╧══╝} \\
\\ 
\texttt{```} \\
\\ 
\texttt{* El borde superior muestra coordenadas x que aumentan hacia la derecha, y el borde derecho coordenadas y que aumentan hacia abajo.} \\
\texttt{* Los siguientes objetos están colocados aleatoriamente en tu cuadrícula: 'W', 'I', 'T', 'C', 'H'.} \\
\\ 
\\ 
\texttt{El otro jugador ve una variación del tablero de juego, donde los objetos están colocados en diferentes ubicaciones aleatorias. No puedes ver el tablero del otro jugador, y él no puede ver el tuyo.} \\
\\ 
\\ 
\texttt{**Objetivo:**} \\
\\ 
\texttt{Ambos jugadores necesitan mover los objetos en sus respectivos fondos de manera que los objetos idénticos terminen en las mismas coordenadas. Tienes que comunicarte con el otro jugador para acordar una configuración objetivo común.} \\
\\ 
\\ 
\texttt{**Reglas:**} \\
\\ 
\texttt{* En cada turno, puedes enviar exactamente uno de los dos comandos siguientes:} \\
\texttt{ 1. `DECIR: <MENSAJE>`: para enviar un mensaje (todo hasta el siguiente salto de línea) al otro jugador. Se lo reenviaré a tu compañero.} \\
\texttt{2. `MOVERR: <OBJETO>, (<X>, <Y>)`: para mover un objeto a una nueva posición, donde `<X>` es la columna y `<Y>` es la fila. Te informaré si tu movimiento fue exitoso o no.} \\
\texttt{* Si no sigues el formato, o envías varios comandos a la vez, tendré que penalizarte.} \\
\texttt{* Si ambos jugadores acumulan más de 12 penalizaciones, ambos pierden el juego.} \\
\texttt{* ¡Es vital que te comuniques con el otro jugador respecto a tu estado objetivo! ¡La *única* forma de transmitir tu estrategia al otro jugador es usando el comando `DECIR: <MENSAJE>`!} \\
\\ 
\\ 
\texttt{**Mover objetos**} \\
\\ 
\texttt{* Solo puedes mover objetos a celdas dentro de los límites de la cuadrícula. La celda de destino debe estar vacía, es decir, solo debe contener el símbolo '◌'.} \\
\texttt{* Si intentas mover un objeto a un lugar que no está vacío, o intentas moverlo fuera de la cuadrícula, tendré que penalizarte. Tendrás otro intento.} \\
\texttt{* ¡Antes de hacer un movimiento, verifica dos veces que el lugar de destino esté vacío, y no contenga ninguna letra, marco o línea!} \\
\\ 
\\ 
\texttt{**Fin del juego**} \\
\\ 
\texttt{Si crees que alcanzaste el objetivo de alinear todos los objetos, puedes pedirle al otro jugador que termine el juego enviando `DECIR: ¿terminado?`. Si el otro jugador te pide que termines el juego, y respondes con `DECIR: ¡terminado!`, el juego terminará.} \\
\\ 
\texttt{Ambos jugadores ganan si el juego termina dentro de 20 rondas, donde una ronda se define como dos jugadores enviando cada uno un comando válido.} \\
\\ 
\texttt{**Puntuación:**} \\
\\ 
\texttt{Cuanto más cerca estén los objetos idénticos en ambos tableros de juego, más puntos obtendrás. Las penalizaciones reducen tus puntos. ¿Puedes batir el récord?} \\
\\ 
\\ 
\texttt{¡Por favor, envía un mensaje al otro jugador para iniciar el juego!} \\
            }
        }
    }
    & \\ \\

    \theutterance \stepcounter{utterance}  
    & \multicolumn{4}{p{0.6\linewidth}}{
        \cellcolor[rgb]{0.9,0.9,0.9}{
            \makecell[{{p{\linewidth}}}]{
                \texttt{\tiny{[P1$\rangle$GM]}}
                \texttt{DECIR: Propongo como objetivo estas coordenadas: W(1,3), I(2,3), T(4,1), H(4,5), C(7,5). Si prefieres otro objetivo, mándame tus coordenadas.} \\
            }
        }
    }
    & & \\ \\

    \theutterance \stepcounter{utterance}  
    & & & \multicolumn{2}{p{0.3\linewidth}}{
        \cellcolor[rgb]{0.9,0.9,0.9}{
            \makecell[{{p{\linewidth}}}]{
                \texttt{\tiny{[GM$|$GM]}}
                \texttt{DECIR: Propongo como objetivo estas coordenadas: W(1,3), I(2,3), T(4,1), H(4,5), C(7,5). Si prefieres otro objetivo, mándame tus coordenadas.} \\
            }
        }
    }
    & & \\ \\

    \theutterance \stepcounter{utterance}  
    & & \multicolumn{4}{p{0.6\linewidth}}{
        \cellcolor[rgb]{0.9,0.9,0.9}{
            \makecell[{{p{\linewidth}}}]{
                \texttt{\tiny{[GM$\rangle$P2]}}
                \texttt{Soy tu maestro de juego, y estás jugando un juego colaborativo con la siguiente cuadrícula como tablero de juego:} \\
\\ 
\texttt{```} \\
\texttt{ 1234567} \\
\texttt{╔══╤══╤═╗} \\
\texttt{║◌◌│W◌│I║ 1} \\
\texttt{╟──┼──┴─╢ 2} \\
\texttt{║◌◌│◌◌◌◌║ 3} \\
\texttt{╟──┘T◌◌◌║ 4} \\
\texttt{║H◌◌◌◌◌◌║ 5} \\
\texttt{║◌◌┌─┬──╢ 6} \\
\texttt{║C◌│◌│◌◌║ 7} \\
\texttt{╚══╧═╧══╝} \\
\\ 
\texttt{```} \\
\\ 
\texttt{* El borde superior muestra coordenadas x que aumentan hacia la derecha, y el borde derecho coordenadas y que aumentan hacia abajo.} \\
\texttt{* Los siguientes objetos están colocados aleatoriamente en tu cuadrícula: 'W', 'I', 'T', 'C', 'H'.} \\
\\ 
\\ 
\texttt{El otro jugador ve una variación del tablero de juego, donde los objetos están colocados en diferentes ubicaciones aleatorias. No puedes ver el tablero del otro jugador, y él no puede ver el tuyo.} \\
\\ 
\\ 
\texttt{**Objetivo:**} \\
\\ 
\texttt{Ambos jugadores necesitan mover los objetos en sus respectivos fondos de manera que los objetos idénticos terminen en las mismas coordenadas. Tienes que comunicarte con el otro jugador para acordar una configuración objetivo común.} \\
\\ 
\\ 
\texttt{**Reglas:**} \\
\\ 
\texttt{* En cada turno, puedes enviar exactamente uno de los dos comandos siguientes:} \\
\texttt{ 1. `DECIR: <MENSAJE>`: para enviar un mensaje (todo hasta el siguiente salto de línea) al otro jugador. Se lo reenviaré a tu compañero.} \\
\texttt{2. `MOVERR: <OBJETO>, (<X>, <Y>)`: para mover un objeto a una nueva posición, donde `<X>` es la columna y `<Y>` es la fila. Te informaré si tu movimiento fue exitoso o no.} \\
\texttt{* Si no sigues el formato, o envías varios comandos a la vez, tendré que penalizarte.} \\
\texttt{* Si ambos jugadores acumulan más de 12 penalizaciones, ambos pierden el juego.} \\
\texttt{* ¡Es vital que te comuniques con el otro jugador respecto a tu estado objetivo! ¡La *única* forma de transmitir tu estrategia al otro jugador es usando el comando `DECIR: <MENSAJE>`!} \\
\\ 
\\ 
\texttt{**Mover objetos**} \\
\\ 
\texttt{* Solo puedes mover objetos a celdas dentro de los límites de la cuadrícula. La celda de destino debe estar vacía, es decir, solo debe contener el símbolo '◌'.} \\
\texttt{* Si intentas mover un objeto a un lugar que no está vacío, o intentas moverlo fuera de la cuadrícula, tendré que penalizarte. Tendrás otro intento.} \\
\texttt{* ¡Antes de hacer un movimiento, verifica dos veces que el lugar de destino esté vacío, y no contenga ninguna letra, marco o línea!} \\
\\ 
\\ 
\texttt{**Fin del juego**} \\
\\ 
\texttt{Si crees que alcanzaste el objetivo de alinear todos los objetos, puedes pedirle al otro jugador que termine el juego enviando `DECIR: ¿terminado?`. Si el otro jugador te pide que termines el juego, y respondes con `DECIR: ¡terminado!`, el juego terminará.} \\
\\ 
\texttt{Ambos jugadores ganan si el juego termina dentro de 20 rondas, donde una ronda se define como dos jugadores enviando cada uno un comando válido.} \\
\\ 
\texttt{**Puntuación:**} \\
\\ 
\texttt{Cuanto más cerca estén los objetos idénticos en ambos tableros de juego, más puntos obtendrás. Las penalizaciones reducen tus puntos. ¿Puedes batir el récord?} \\
\\ 
\texttt{El otro jugador inició el juego enviando este mensaje:} \\
\\ 
\texttt{"Propongo como objetivo estas coordenadas: W(1,3), I(2,3), T(4,1), H(4,5), C(7,5). Si prefieres otro objetivo, mándame tus coordenadas."} \\
\\ 
\texttt{¿Cuál es tu primer comando?} \\
            }
        }
    }
    & \\ \\

    \theutterance \stepcounter{utterance}  
    & & & \multicolumn{4}{p{0.6\linewidth}}{
        \cellcolor[rgb]{0.9,0.9,0.9}{
            \makecell[{{p{\linewidth}}}]{
                \texttt{\tiny{[GM$\langle$P2]}}
                \texttt{DECIR: Acepto tu objetivo: W(1,3), I(2,3), T(4,1), H(4,5), C(7,5).} \\
            }
        }
    }
     \\ \\

    \theutterance \stepcounter{utterance}  
    & & & \multicolumn{2}{p{0.3\linewidth}}{
        \cellcolor[rgb]{0.9,0.9,0.9}{
            \makecell[{{p{\linewidth}}}]{
                \texttt{\tiny{[GM$|$GM]}}
                \texttt{DECIR: Acepto tu objetivo: W(1,3), I(2,3), T(4,1), H(4,5), C(7,5).} \\
            }
        }
    }
    & & \\ \\

    \theutterance \stepcounter{utterance}  
    & & \multicolumn{4}{p{0.6\linewidth}}{
        \cellcolor[rgb]{0.9,0.9,0.9}{
            \makecell[{{p{\linewidth}}}]{
                \texttt{\tiny{[P1$\langle$GM]}}
                \texttt{Tu mensaje ha sido transmitido al otro jugador.} \\
\\ 
\\ 
\texttt{Actualmente estás jugando la ronda 1 de un máximo de 20.} \\
\texttt{Han acumulado colectivamente 0 de 12 penalizaciones.} \\
\texttt{El otro jugador envió este mensaje:} \\
\\ 
\\ 
\texttt{"Acepto tu objetivo: W(1,3), I(2,3), T(4,1), H(4,5), C(7,5)."} \\
\\ 
\\ 
\texttt{¿Cuál es tu próximo comando?} \\
            }
        }
    }
    & \\ \\

    \theutterance \stepcounter{utterance}  
    & \multicolumn{4}{p{0.6\linewidth}}{
        \cellcolor[rgb]{0.9,0.9,0.9}{
            \makecell[{{p{\linewidth}}}]{
                \texttt{\tiny{[P1$\rangle$GM]}}
                \texttt{DECIR: Mi tablero ya coincide con el objetivo acordado; no necesito mover nada. Cuando el tuyo también coincida, envíame: ¿terminado?} \\
            }
        }
    }
    & & \\ \\

    \theutterance \stepcounter{utterance}  
    & & & \multicolumn{2}{p{0.3\linewidth}}{
        \cellcolor[rgb]{0.9,0.9,0.9}{
            \makecell[{{p{\linewidth}}}]{
                \texttt{\tiny{[GM$|$GM]}}
                \texttt{DECIR: Mi tablero ya coincide con el objetivo acordado; no necesito mover nada. Cuando el tuyo también coincida, envíame: ¿terminado?} \\
            }
        }
    }
    & & \\ \\

    \theutterance \stepcounter{utterance}  
    & & \multicolumn{4}{p{0.6\linewidth}}{
        \cellcolor[rgb]{0.9,0.9,0.9}{
            \makecell[{{p{\linewidth}}}]{
                \texttt{\tiny{[GM$\rangle$P2]}}
                \texttt{Tu mensaje ha sido transmitido al otro jugador.} \\
\\ 
\\ 
\texttt{Actualmente estás jugando la ronda 2 de un máximo de 20.} \\
\texttt{Han acumulado colectivamente 0 de 12 penalizaciones.} \\
\texttt{El otro jugador envió este mensaje:} \\
\\ 
\\ 
\texttt{"Mi tablero ya coincide con el objetivo acordado; no necesito mover nada. Cuando el tuyo también coincida, envíame: ¿terminado?"} \\
\\ 
\\ 
\texttt{¿Cuál es tu próximo comando?} \\
            }
        }
    }
    & \\ \\

    \theutterance \stepcounter{utterance}  
    & & & \multicolumn{4}{p{0.6\linewidth}}{
        \cellcolor[rgb]{0.9,0.9,0.9}{
            \makecell[{{p{\linewidth}}}]{
                \texttt{\tiny{[GM$\langle$P2]}}
                \texttt{MOVERR: H, (4, 5)} \\
            }
        }
    }
     \\ \\

    \theutterance \stepcounter{utterance}  
    & & & \multicolumn{2}{p{0.3\linewidth}}{
        \cellcolor[rgb]{0.9,0.9,0.9}{
            \makecell[{{p{\linewidth}}}]{
                \texttt{\tiny{[GM$|$GM]}}
                \texttt{MOVERR: H, (4, 5)} \\
            }
        }
    }
    & & \\ \\

    \theutterance \stepcounter{utterance}  
    & & & \multicolumn{2}{p{0.3\linewidth}}{
        \cellcolor[rgb]{0.9,0.9,0.9}{
            \makecell[{{p{\linewidth}}}]{
                \texttt{\tiny{[GM$|$GM]}}
                \texttt{Invalid response format} \\
            }
        }
    }
    & & \\ \\

    \theutterance \stepcounter{utterance}  
    & & \multicolumn{4}{p{0.6\linewidth}}{
        \cellcolor[rgb]{0.9,0.9,0.9}{
            \makecell[{{p{\linewidth}}}]{
                \texttt{\tiny{[GM$\rangle$P2]}}
                \texttt{Penalización: ¡Tu mensaje no tiene el formato esperado!} \\
\\ 
\texttt{Asegúrate de que tu respuesta solo contenga `DECIR: <MENSAJE>` o `MOVERR: <OBJETO>, (<X>, <Y>)`, ¡y nada más!} \\
\\ 
\texttt{Han acumulado colectivamente 1 de 12 penalizaciones.} \\
\texttt{¡Por favor, inténtalo de nuevo!} \\
            }
        }
    }
    & \\ \\

    \theutterance \stepcounter{utterance}  
    & & & \multicolumn{4}{p{0.6\linewidth}}{
        \cellcolor[rgb]{0.9,0.9,0.9}{
            \makecell[{{p{\linewidth}}}]{
                \texttt{\tiny{[GM$\langle$P2]}}
                \texttt{MOVERR: W, (1, 3)} \\
            }
        }
    }
     \\ \\

    \theutterance \stepcounter{utterance}  
    & & & \multicolumn{2}{p{0.3\linewidth}}{
        \cellcolor[rgb]{0.9,0.9,0.9}{
            \makecell[{{p{\linewidth}}}]{
                \texttt{\tiny{[GM$|$GM]}}
                \texttt{MOVERR: W, (1, 3)} \\
            }
        }
    }
    & & \\ \\

    \theutterance \stepcounter{utterance}  
    & & & \multicolumn{2}{p{0.3\linewidth}}{
        \cellcolor[rgb]{0.9,0.9,0.9}{
            \makecell[{{p{\linewidth}}}]{
                \texttt{\tiny{[GM$|$GM]}}
                \texttt{Invalid response format} \\
            }
        }
    }
    & & \\ \\

    \theutterance \stepcounter{utterance}  
    & & \multicolumn{4}{p{0.6\linewidth}}{
        \cellcolor[rgb]{0.9,0.9,0.9}{
            \makecell[{{p{\linewidth}}}]{
                \texttt{\tiny{[GM$\rangle$P2]}}
                \texttt{Penalización: ¡Tu mensaje no tiene el formato esperado!} \\
\\ 
\texttt{Asegúrate de que tu respuesta solo contenga `DECIR: <MENSAJE>` o `MOVERR: <OBJETO>, (<X>, <Y>)`, ¡y nada más!} \\
\\ 
\texttt{Han acumulado colectivamente 2 de 12 penalizaciones.} \\
\texttt{¡Por favor, inténtalo de nuevo!} \\
            }
        }
    }
    & \\ \\

    \theutterance \stepcounter{utterance}  
    & & & \multicolumn{4}{p{0.6\linewidth}}{
        \cellcolor[rgb]{0.9,0.9,0.9}{
            \makecell[{{p{\linewidth}}}]{
                \texttt{\tiny{[GM$\langle$P2]}}
                \texttt{MOVER: W, (1, 3)} \\
            }
        }
    }
     \\ \\

    \theutterance \stepcounter{utterance}  
    & & & \multicolumn{2}{p{0.3\linewidth}}{
        \cellcolor[rgb]{0.9,0.9,0.9}{
            \makecell[{{p{\linewidth}}}]{
                \texttt{\tiny{[GM$|$GM]}}
                \texttt{MOVER: W, (1, 3)} \\
            }
        }
    }
    & & \\ \\

    \theutterance \stepcounter{utterance}  
    & & & \multicolumn{2}{p{0.3\linewidth}}{
        \cellcolor[rgb]{0.9,0.9,0.9}{
            \makecell[{{p{\linewidth}}}]{
                \texttt{\tiny{[GM$|$GM]}}
                \texttt{Invalid response format} \\
            }
        }
    }
    & & \\ \\

    \theutterance \stepcounter{utterance}  
    & & \multicolumn{4}{p{0.6\linewidth}}{
        \cellcolor[rgb]{0.9,0.9,0.9}{
            \makecell[{{p{\linewidth}}}]{
                \texttt{\tiny{[GM$\rangle$P2]}}
                \texttt{Penalización: ¡Tu mensaje no tiene el formato esperado!} \\
\\ 
\texttt{Asegúrate de que tu respuesta solo contenga `DECIR: <MENSAJE>` o `MOVERR: <OBJETO>, (<X>, <Y>)`, ¡y nada más!} \\
\\ 
\texttt{Han acumulado colectivamente 3 de 12 penalizaciones.} \\
\texttt{¡Por favor, inténtalo de nuevo!} \\
            }
        }
    }
    & \\ \\

    \theutterance \stepcounter{utterance}  
    & & & \multicolumn{4}{p{0.6\linewidth}}{
        \cellcolor[rgb]{0.9,0.9,0.9}{
            \makecell[{{p{\linewidth}}}]{
                \texttt{\tiny{[GM$\langle$P2]}}
                \texttt{MOVE: W, (1, 3)} \\
            }
        }
    }
     \\ \\

    \theutterance \stepcounter{utterance}  
    & & & \multicolumn{2}{p{0.3\linewidth}}{
        \cellcolor[rgb]{0.9,0.9,0.9}{
            \makecell[{{p{\linewidth}}}]{
                \texttt{\tiny{[GM$|$GM]}}
                \texttt{MOVE: W, (1, 3)} \\
            }
        }
    }
    & & \\ \\

    \theutterance \stepcounter{utterance}  
    & & & \multicolumn{2}{p{0.3\linewidth}}{
        \cellcolor[rgb]{0.9,0.9,0.9}{
            \makecell[{{p{\linewidth}}}]{
                \texttt{\tiny{[GM$|$GM]}}
                \texttt{Invalid response format} \\
            }
        }
    }
    & & \\ \\

    \theutterance \stepcounter{utterance}  
    & & \multicolumn{4}{p{0.6\linewidth}}{
        \cellcolor[rgb]{0.9,0.9,0.9}{
            \makecell[{{p{\linewidth}}}]{
                \texttt{\tiny{[GM$\rangle$P2]}}
                \texttt{Penalización: ¡Tu mensaje no tiene el formato esperado!} \\
\\ 
\texttt{Asegúrate de que tu respuesta solo contenga `DECIR: <MENSAJE>` o `MOVERR: <OBJETO>, (<X>, <Y>)`, ¡y nada más!} \\
\\ 
\texttt{Han acumulado colectivamente 4 de 12 penalizaciones.} \\
\texttt{¡Por favor, inténtalo de nuevo!} \\
            }
        }
    }
    & \\ \\

    \theutterance \stepcounter{utterance}  
    & & & \multicolumn{4}{p{0.6\linewidth}}{
        \cellcolor[rgb]{0.9,0.9,0.9}{
            \makecell[{{p{\linewidth}}}]{
                \texttt{\tiny{[GM$\langle$P2]}}
                \texttt{MOVERR:W,(1,3)} \\
            }
        }
    }
     \\ \\

    \theutterance \stepcounter{utterance}  
    & & & \multicolumn{2}{p{0.3\linewidth}}{
        \cellcolor[rgb]{0.9,0.9,0.9}{
            \makecell[{{p{\linewidth}}}]{
                \texttt{\tiny{[GM$|$GM]}}
                \texttt{MOVERR:W,(1,3)} \\
            }
        }
    }
    & & \\ \\

    \theutterance \stepcounter{utterance}  
    & & & \multicolumn{2}{p{0.3\linewidth}}{
        \cellcolor[rgb]{0.9,0.9,0.9}{
            \makecell[{{p{\linewidth}}}]{
                \texttt{\tiny{[GM$|$GM]}}
                \texttt{Invalid response format} \\
            }
        }
    }
    & & \\ \\

    \theutterance \stepcounter{utterance}  
    & & \multicolumn{4}{p{0.6\linewidth}}{
        \cellcolor[rgb]{0.9,0.9,0.9}{
            \makecell[{{p{\linewidth}}}]{
                \texttt{\tiny{[GM$\rangle$P2]}}
                \texttt{Penalización: ¡Tu mensaje no tiene el formato esperado!} \\
\\ 
\texttt{Asegúrate de que tu respuesta solo contenga `DECIR: <MENSAJE>` o `MOVERR: <OBJETO>, (<X>, <Y>)`, ¡y nada más!} \\
\\ 
\texttt{Han acumulado colectivamente 5 de 12 penalizaciones.} \\
\texttt{¡Por favor, inténtalo de nuevo!} \\
            }
        }
    }
    & \\ \\

    \theutterance \stepcounter{utterance}  
    & & & \multicolumn{4}{p{0.6\linewidth}}{
        \cellcolor[rgb]{0.9,0.9,0.9}{
            \makecell[{{p{\linewidth}}}]{
                \texttt{\tiny{[GM$\langle$P2]}}
                \texttt{MOVER: W, (1,3)} \\
            }
        }
    }
     \\ \\

    \theutterance \stepcounter{utterance}  
    & & & \multicolumn{2}{p{0.3\linewidth}}{
        \cellcolor[rgb]{0.9,0.9,0.9}{
            \makecell[{{p{\linewidth}}}]{
                \texttt{\tiny{[GM$|$GM]}}
                \texttt{MOVER: W, (1,3)} \\
            }
        }
    }
    & & \\ \\

    \theutterance \stepcounter{utterance}  
    & & & \multicolumn{2}{p{0.3\linewidth}}{
        \cellcolor[rgb]{0.9,0.9,0.9}{
            \makecell[{{p{\linewidth}}}]{
                \texttt{\tiny{[GM$|$GM]}}
                \texttt{Invalid response format} \\
            }
        }
    }
    & & \\ \\

    \theutterance \stepcounter{utterance}  
    & & \multicolumn{4}{p{0.6\linewidth}}{
        \cellcolor[rgb]{0.9,0.9,0.9}{
            \makecell[{{p{\linewidth}}}]{
                \texttt{\tiny{[GM$\rangle$P2]}}
                \texttt{Penalización: ¡Tu mensaje no tiene el formato esperado!} \\
\\ 
\texttt{Asegúrate de que tu respuesta solo contenga `DECIR: <MENSAJE>` o `MOVERR: <OBJETO>, (<X>, <Y>)`, ¡y nada más!} \\
\\ 
\texttt{Han acumulado colectivamente 6 de 12 penalizaciones.} \\
\texttt{¡Por favor, inténtalo de nuevo!} \\
            }
        }
    }
    & \\ \\

    \theutterance \stepcounter{utterance}  
    & & & \multicolumn{4}{p{0.6\linewidth}}{
        \cellcolor[rgb]{0.9,0.9,0.9}{
            \makecell[{{p{\linewidth}}}]{
                \texttt{\tiny{[GM$\langle$P2]}}
                \texttt{DECIR: Estoy teniendo problemas con el formato de MOVERR. Mi objetivo sigue siendo W(1,3), I(2,3), T(4,1), H(4,5), C(7,5).} \\
            }
        }
    }
     \\ \\

    \theutterance \stepcounter{utterance}  
    & & & \multicolumn{2}{p{0.3\linewidth}}{
        \cellcolor[rgb]{0.9,0.9,0.9}{
            \makecell[{{p{\linewidth}}}]{
                \texttt{\tiny{[GM$|$GM]}}
                \texttt{DECIR: Estoy teniendo problemas con el formato de MOVERR. Mi objetivo sigue siendo W(1,3), I(2,3), T(4,1), H(4,5), C(7,5).} \\
            }
        }
    }
    & & \\ \\

    \theutterance \stepcounter{utterance}  
    & & \multicolumn{4}{p{0.6\linewidth}}{
        \cellcolor[rgb]{0.9,0.9,0.9}{
            \makecell[{{p{\linewidth}}}]{
                \texttt{\tiny{[P1$\langle$GM]}}
                \texttt{Tu mensaje ha sido transmitido al otro jugador.} \\
\\ 
\\ 
\texttt{Actualmente estás jugando la ronda 2 de un máximo de 20.} \\
\texttt{Han acumulado colectivamente 6 de 12 penalizaciones.} \\
\texttt{El otro jugador envió este mensaje:} \\
\\ 
\\ 
\texttt{"Estoy teniendo problemas con el formato de MOVERR. Mi objetivo sigue siendo W(1,3), I(2,3), T(4,1), H(4,5), C(7,5)."} \\
\\ 
\\ 
\texttt{¿Cuál es tu próximo comando?} \\
            }
        }
    }
    & \\ \\

    \theutterance \stepcounter{utterance}  
    & \multicolumn{4}{p{0.6\linewidth}}{
        \cellcolor[rgb]{0.9,0.9,0.9}{
            \makecell[{{p{\linewidth}}}]{
                \texttt{\tiny{[P1$\rangle$GM]}}
                \texttt{DECIR: Usa exactamente este formato, una sola línea y un solo comando por turno: MOVERR: W, (1, 3)  ejemplo. Cambia letra y coordenadas según necesites. No pongas texto extra en ese turno. Cuando termines tu tablero, mándame: ¿terminado?} \\
            }
        }
    }
    & & \\ \\

    \theutterance \stepcounter{utterance}  
    & & & \multicolumn{2}{p{0.3\linewidth}}{
        \cellcolor[rgb]{0.9,0.9,0.9}{
            \makecell[{{p{\linewidth}}}]{
                \texttt{\tiny{[GM$|$GM]}}
                \texttt{DECIR: Usa exactamente este formato, una sola línea y un solo comando por turno: MOVERR: W, (1, 3)  ejemplo. Cambia letra y coordenadas según necesites. No pongas texto extra en ese turno. Cuando termines tu tablero, mándame: ¿terminado?} \\
            }
        }
    }
    & & \\ \\

    \theutterance \stepcounter{utterance}  
    & & \multicolumn{4}{p{0.6\linewidth}}{
        \cellcolor[rgb]{0.9,0.9,0.9}{
            \makecell[{{p{\linewidth}}}]{
                \texttt{\tiny{[GM$\rangle$P2]}}
                \texttt{Tu mensaje ha sido transmitido al otro jugador.} \\
\\ 
\\ 
\texttt{Actualmente estás jugando la ronda 3 de un máximo de 20.} \\
\texttt{Han acumulado colectivamente 6 de 12 penalizaciones.} \\
\texttt{El otro jugador envió este mensaje:} \\
\\ 
\\ 
\texttt{"Usa exactamente este formato, una sola línea y un solo comando por turno: MOVERR: W, (1, 3)  ejemplo. Cambia letra y coordenadas según necesites. No pongas texto extra en ese turno. Cuando termines tu tablero, mándame: ¿terminado?"} \\
\\ 
\\ 
\texttt{¿Cuál es tu próximo comando?} \\
            }
        }
    }
    & \\ \\

    \theutterance \stepcounter{utterance}  
    & & & \multicolumn{4}{p{0.6\linewidth}}{
        \cellcolor[rgb]{0.9,0.9,0.9}{
            \makecell[{{p{\linewidth}}}]{
                \texttt{\tiny{[GM$\langle$P2]}}
                \texttt{MOVERR: W, (1, 3)} \\
            }
        }
    }
     \\ \\

    \theutterance \stepcounter{utterance}  
    & & & \multicolumn{2}{p{0.3\linewidth}}{
        \cellcolor[rgb]{0.9,0.9,0.9}{
            \makecell[{{p{\linewidth}}}]{
                \texttt{\tiny{[GM$|$GM]}}
                \texttt{MOVERR: W, (1, 3)} \\
            }
        }
    }
    & & \\ \\

    \theutterance \stepcounter{utterance}  
    & & & \multicolumn{2}{p{0.3\linewidth}}{
        \cellcolor[rgb]{0.9,0.9,0.9}{
            \makecell[{{p{\linewidth}}}]{
                \texttt{\tiny{[GM$|$GM]}}
                \texttt{Invalid response format} \\
            }
        }
    }
    & & \\ \\

    \theutterance \stepcounter{utterance}  
    & & \multicolumn{4}{p{0.6\linewidth}}{
        \cellcolor[rgb]{0.9,0.9,0.9}{
            \makecell[{{p{\linewidth}}}]{
                \texttt{\tiny{[GM$\rangle$P2]}}
                \texttt{Penalización: ¡Tu mensaje no tiene el formato esperado!} \\
\\ 
\texttt{Asegúrate de que tu respuesta solo contenga `DECIR: <MENSAJE>` o `MOVERR: <OBJETO>, (<X>, <Y>)`, ¡y nada más!} \\
\\ 
\texttt{Han acumulado colectivamente 7 de 12 penalizaciones.} \\
\texttt{¡Por favor, inténtalo de nuevo!} \\
            }
        }
    }
    & \\ \\

    \theutterance \stepcounter{utterance}  
    & & & \multicolumn{4}{p{0.6\linewidth}}{
        \cellcolor[rgb]{0.9,0.9,0.9}{
            \makecell[{{p{\linewidth}}}]{
                \texttt{\tiny{[GM$\langle$P2]}}
                \texttt{DECIR: El comando de mover sigue siendo rechazado incluso copiando exactamente tu ejemplo. De momento evito más penalizaciones; si sabes otra sintaxis válida, dímela.} \\
            }
        }
    }
     \\ \\

    \theutterance \stepcounter{utterance}  
    & & & \multicolumn{2}{p{0.3\linewidth}}{
        \cellcolor[rgb]{0.9,0.9,0.9}{
            \makecell[{{p{\linewidth}}}]{
                \texttt{\tiny{[GM$|$GM]}}
                \texttt{DECIR: El comando de mover sigue siendo rechazado incluso copiando exactamente tu ejemplo. De momento evito más penalizaciones; si sabes otra sintaxis válida, dímela.} \\
            }
        }
    }
    & & \\ \\

    \theutterance \stepcounter{utterance}  
    & & \multicolumn{4}{p{0.6\linewidth}}{
        \cellcolor[rgb]{0.9,0.9,0.9}{
            \makecell[{{p{\linewidth}}}]{
                \texttt{\tiny{[P1$\langle$GM]}}
                \texttt{Tu mensaje ha sido transmitido al otro jugador.} \\
\\ 
\\ 
\texttt{Actualmente estás jugando la ronda 3 de un máximo de 20.} \\
\texttt{Han acumulado colectivamente 7 de 12 penalizaciones.} \\
\texttt{El otro jugador envió este mensaje:} \\
\\ 
\\ 
\texttt{"El comando de mover sigue siendo rechazado incluso copiando exactamente tu ejemplo. De momento evito más penalizaciones; si sabes otra sintaxis válida, dímela."} \\
\\ 
\\ 
\texttt{¿Cuál es tu próximo comando?} \\
            }
        }
    }
    & \\ \\

    \theutterance \stepcounter{utterance}  
    & \multicolumn{4}{p{0.6\linewidth}}{
        \cellcolor[rgb]{0.9,0.9,0.9}{
            \makecell[{{p{\linewidth}}}]{
                \texttt{\tiny{[P1$\rangle$GM]}}
                \texttt{DECIR: Prueba esta variante exacta, sin espacio tras la letra ni dentro de coordenadas: MOVERR: W,(1,3). Si funciona, usa siempre ese patrón: MOVERR: LETRA,(x,y).} \\
            }
        }
    }
    & & \\ \\

    \theutterance \stepcounter{utterance}  
    & & & \multicolumn{2}{p{0.3\linewidth}}{
        \cellcolor[rgb]{0.9,0.9,0.9}{
            \makecell[{{p{\linewidth}}}]{
                \texttt{\tiny{[GM$|$GM]}}
                \texttt{DECIR: Prueba esta variante exacta, sin espacio tras la letra ni dentro de coordenadas: MOVERR: W,(1,3). Si funciona, usa siempre ese patrón: MOVERR: LETRA,(x,y).} \\
            }
        }
    }
    & & \\ \\

    \theutterance \stepcounter{utterance}  
    & & \multicolumn{4}{p{0.6\linewidth}}{
        \cellcolor[rgb]{0.9,0.9,0.9}{
            \makecell[{{p{\linewidth}}}]{
                \texttt{\tiny{[GM$\rangle$P2]}}
                \texttt{Tu mensaje ha sido transmitido al otro jugador.} \\
\\ 
\\ 
\texttt{Actualmente estás jugando la ronda 4 de un máximo de 20.} \\
\texttt{Han acumulado colectivamente 7 de 12 penalizaciones.} \\
\texttt{El otro jugador envió este mensaje:} \\
\\ 
\\ 
\texttt{"Prueba esta variante exacta, sin espacio tras la letra ni dentro de coordenadas: MOVERR: W,(1,3). Si funciona, usa siempre ese patrón: MOVERR: LETRA,(x,y)."} \\
\\ 
\\ 
\texttt{¿Cuál es tu próximo comando?} \\
            }
        }
    }
    & \\ \\

    \theutterance \stepcounter{utterance}  
    & & & \multicolumn{4}{p{0.6\linewidth}}{
        \cellcolor[rgb]{0.9,0.9,0.9}{
            \makecell[{{p{\linewidth}}}]{
                \texttt{\tiny{[GM$\langle$P2]}}
                \texttt{MOVERR: W,(1,3)} \\
            }
        }
    }
     \\ \\

    \theutterance \stepcounter{utterance}  
    & & & \multicolumn{2}{p{0.3\linewidth}}{
        \cellcolor[rgb]{0.9,0.9,0.9}{
            \makecell[{{p{\linewidth}}}]{
                \texttt{\tiny{[GM$|$GM]}}
                \texttt{MOVERR: W,(1,3)} \\
            }
        }
    }
    & & \\ \\

    \theutterance \stepcounter{utterance}  
    & & & \multicolumn{2}{p{0.3\linewidth}}{
        \cellcolor[rgb]{0.9,0.9,0.9}{
            \makecell[{{p{\linewidth}}}]{
                \texttt{\tiny{[GM$|$GM]}}
                \texttt{Invalid response format} \\
            }
        }
    }
    & & \\ \\

    \theutterance \stepcounter{utterance}  
    & & \multicolumn{4}{p{0.6\linewidth}}{
        \cellcolor[rgb]{0.9,0.9,0.9}{
            \makecell[{{p{\linewidth}}}]{
                \texttt{\tiny{[GM$\rangle$P2]}}
                \texttt{Penalización: ¡Tu mensaje no tiene el formato esperado!} \\
\\ 
\texttt{Asegúrate de que tu respuesta solo contenga `DECIR: <MENSAJE>` o `MOVERR: <OBJETO>, (<X>, <Y>)`, ¡y nada más!} \\
\\ 
\texttt{Han acumulado colectivamente 8 de 12 penalizaciones.} \\
\texttt{¡Por favor, inténtalo de nuevo!} \\
            }
        }
    }
    & \\ \\

    \theutterance \stepcounter{utterance}  
    & & & \multicolumn{4}{p{0.6\linewidth}}{
        \cellcolor[rgb]{0.9,0.9,0.9}{
            \makecell[{{p{\linewidth}}}]{
                \texttt{\tiny{[GM$\langle$P2]}}
                \texttt{DECIR: No puedo ejecutar movimientos porque el sistema rechaza todos los formatos de MOVERR. Para evitar más penalizaciones, propongo terminar si te parece aceptable la puntuación actual: ¿terminado?} \\
            }
        }
    }
     \\ \\

    \theutterance \stepcounter{utterance}  
    & & & \multicolumn{2}{p{0.3\linewidth}}{
        \cellcolor[rgb]{0.9,0.9,0.9}{
            \makecell[{{p{\linewidth}}}]{
                \texttt{\tiny{[GM$|$GM]}}
                \texttt{DECIR: No puedo ejecutar movimientos porque el sistema rechaza todos los formatos de MOVERR. Para evitar más penalizaciones, propongo terminar si te parece aceptable la puntuación actual: ¿terminado?} \\
            }
        }
    }
    & & \\ \\

    \theutterance \stepcounter{utterance}  
    & & \multicolumn{4}{p{0.6\linewidth}}{
        \cellcolor[rgb]{0.9,0.9,0.9}{
            \makecell[{{p{\linewidth}}}]{
                \texttt{\tiny{[P1$\langle$GM]}}
                \texttt{Tu mensaje ha sido transmitido al otro jugador.} \\
\\ 
\\ 
\texttt{Actualmente estás jugando la ronda 4 de un máximo de 20.} \\
\texttt{Han acumulado colectivamente 8 de 12 penalizaciones.} \\
\texttt{El otro jugador envió este mensaje:} \\
\\ 
\\ 
\texttt{"No puedo ejecutar movimientos porque el sistema rechaza todos los formatos de MOVERR. Para evitar más penalizaciones, propongo terminar si te parece aceptable la puntuación actual: ¿terminado?"} \\
\\ 
\\ 
\texttt{¿Cuál es tu próximo comando?} \\
            }
        }
    }
    & \\ \\

    \theutterance \stepcounter{utterance}  
    & \multicolumn{4}{p{0.6\linewidth}}{
        \cellcolor[rgb]{0.9,0.9,0.9}{
            \makecell[{{p{\linewidth}}}]{
                \texttt{\tiny{[P1$\rangle$GM]}}
                \texttt{DECIR: ¡terminado!} \\
            }
        }
    }
    & & \\ \\

    \theutterance \stepcounter{utterance}  
    & & & \multicolumn{2}{p{0.3\linewidth}}{
        \cellcolor[rgb]{0.9,0.9,0.9}{
            \makecell[{{p{\linewidth}}}]{
                \texttt{\tiny{[GM$|$GM]}}
                \texttt{DECIR: ¡terminado!} \\
            }
        }
    }
    & & \\ \\

    \theutterance \stepcounter{utterance}  
    & & & \multicolumn{2}{p{0.3\linewidth}}{
        \cellcolor[rgb]{0.9,0.9,0.9}{
            \makecell[{{p{\linewidth}}}]{
                \texttt{\tiny{[GM$|$GM]}}
                \texttt{true} \\
            }
        }
    }
    & & \\ \\

    \theutterance \stepcounter{utterance}  
    & & & \multicolumn{2}{p{0.3\linewidth}}{
        \cellcolor[rgb]{0.9,0.9,0.9}{
            \makecell[{{p{\linewidth}}}]{
                \texttt{\tiny{[GM$|$GM]}}
                \texttt{Initial states:} \\
\texttt{```} \\
\texttt{Player 1:} \\
\texttt{ 1234567} \\
\texttt{╔══╤══╤═╗} \\
\texttt{║◌◌│T◌│◌║ 1} \\
\texttt{╟──┼──┴─╢ 2} \\
\texttt{║WI│◌◌◌◌║ 3} \\
\texttt{╟──┘◌◌◌◌║ 4} \\
\texttt{║◌◌◌H◌◌C║ 5} \\
\texttt{║◌◌┌─┬──╢ 6} \\
\texttt{║◌◌│◌│◌◌║ 7} \\
\texttt{╚══╧═╧══╝} \\
\\ 
\\ 
\texttt{Player 2:} \\
\texttt{ 1234567} \\
\texttt{╔══╤══╤═╗} \\
\texttt{║◌◌│W◌│I║ 1} \\
\texttt{╟──┼──┴─╢ 2} \\
\texttt{║◌◌│◌◌◌◌║ 3} \\
\texttt{╟──┘T◌◌◌║ 4} \\
\texttt{║H◌◌◌◌◌◌║ 5} \\
\texttt{║◌◌┌─┬──╢ 6} \\
\texttt{║C◌│◌│◌◌║ 7} \\
\texttt{╚══╧═╧══╝} \\
\\ 
\texttt{```} \\
            }
        }
    }
    & & \\ \\

    \theutterance \stepcounter{utterance}  
    & & & \multicolumn{2}{p{0.3\linewidth}}{
        \cellcolor[rgb]{0.9,0.9,0.9}{
            \makecell[{{p{\linewidth}}}]{
                \texttt{\tiny{[GM$|$GM]}}
                \texttt{End states:} \\
\texttt{```} \\
\texttt{Player 1:} \\
\texttt{ 1234567} \\
\texttt{╔══╤══╤═╗} \\
\texttt{║◌◌│T◌│◌║ 1} \\
\texttt{╟──┼──┴─╢ 2} \\
\texttt{║WI│◌◌◌◌║ 3} \\
\texttt{╟──┘◌◌◌◌║ 4} \\
\texttt{║◌◌◌H◌◌C║ 5} \\
\texttt{║◌◌┌─┬──╢ 6} \\
\texttt{║◌◌│◌│◌◌║ 7} \\
\texttt{╚══╧═╧══╝} \\
\\ 
\\ 
\texttt{Player 2:} \\
\texttt{ 1234567} \\
\texttt{╔══╤══╤═╗} \\
\texttt{║◌◌│W◌│I║ 1} \\
\texttt{╟──┼──┴─╢ 2} \\
\texttt{║◌◌│◌◌◌◌║ 3} \\
\texttt{╟──┘T◌◌◌║ 4} \\
\texttt{║H◌◌◌◌◌◌║ 5} \\
\texttt{║◌◌┌─┬──╢ 6} \\
\texttt{║C◌│◌│◌◌║ 7} \\
\texttt{╚══╧═╧══╝} \\
\\ 
\texttt{```} \\
            }
        }
    }
    & & \\ \\

    \theutterance \stepcounter{utterance}  
    & & & \multicolumn{2}{p{0.3\linewidth}}{
        \cellcolor[rgb]{0.9,0.9,0.9}{
            \makecell[{{p{\linewidth}}}]{
                \texttt{\tiny{[GM$|$GM]}}
                \texttt{* success: True} \\
\texttt{* lose: True} \\
\texttt{* aborted: False} \\
\texttt{{-}{-}{-}{-}{-}{-}{-}} \\
\texttt{* Shifts: 0.00} \\
\texttt{* Max Shifts: 8.00} \\
\texttt{* Min Shifts: 4.00} \\
\texttt{* End Distance Sum: 21.32} \\
\texttt{* Init Distance Sum: 21.32} \\
\texttt{* Expected Distance Sum: 20.95} \\
\texttt{* Penalties: 8.00} \\
\texttt{* Max Penalties: 12.00} \\
\texttt{* Rounds: 4.00} \\
\texttt{* Max Rounds: 20.00} \\
\texttt{* Object Count: 5.00} \\
            }
        }
    }
    & & \\ \\

\end{supertabular}
}

\end{document}
